\color{uniblue}{\section{Техническое обслуживание и ремонт}}
\color{black}

\color{uniblue}{\subsection{\raggedright{Техническое обслуживание устройства}}}
\color{black}
\nopagebreak
\vspace{3mm}
\par
\sloppy

\par
Процесс  технического  обслуживания  устройств  релейной  защиты  регламентирован  приказом Минэнерго  России  «Об  утверждении  Правил  технического  обслуживания  устройств  и  комплексов релейной  защиты  и  автоматики  и  внесении  изменений  в  требования  к  обеспечению  надежности электроэнергетических  систем,  надежности  и  безопасности  объектов  электроэнергетики  и 
энергопринимающих установок». 
Техническое обслуживание устройств ЮНИТ-М300 должно выполняться в соответствии с документом \color{red}ЮТКБ.656112.608  ТО  –  «Микропроцессорное  устройство  релейной  защиты  и  автоматики  серии «ЮНИТ-М300»\color{black}. Методические указания по выполнению технического обслуживания». 

{
\color{uniblue}{\subsection{\raggedright{Текущий ремонт}}}
\color{black}
\nopagebreak
\vspace{3mm}
}

Устройство  ЮНИТ-М300,  как  было  описано  выше,  выполнено  с  использованием  инновационных микропроцессорных  технологий,  поэтому  необходимость  в  его  текущем  ремонте  отсутствует. 
Работоспособность  и  безотказность  работы  устройства  обеспечивается  функционированием  системы самодиагностики. 
Углубленная самодиагностика аппаратных средств, а также состояние программной памяти, памяти данных  и  буферной  памяти  производится  при  включении  устройства.  Поэтому  при  возникновении 
сообщений о критических или некритических ошибках от системы самодиагностики, даже при возможности устранения  их  своими  силами,  следует  обратиться  на  предприятие-изготовитель для  консультации или проведения регламентных работ. 
Поводом  для  обращения  на  предприятие-изготовитель  являются  также  признаки  нехарактерного нагрева или перегрева отдельных элементов, участков или частей устройства. 
Вышедшие  из  строя  устройства  должны  проходить  ремонт  на  предприятии-изготовителе, обеспечивающем  гарантийное  и  послегарантийное  обслуживание,  адрес  которого  указан  в  паспорте 
устройства.  В  качестве  ЗИП,  если  это  предусмотрено  договором  поставки,  могут  предоставляться устройства  ЮНИТ-М300  с  необходимыми  конфигурациями,  позволяющие  оперативно  произвести  замену вышедших из строя устройств.

